\documentclass[10pt]{article}

\usepackage[letterpaper,hmargin=0.82in,vmargin=0.72in]{geometry}
\usepackage[T1]{fontenc}
\usepackage{amsmath,amssymb,amsthm}
\usepackage{newtxtext,newtxmath}
\usepackage{microtype}
\usepackage{xcolor}
\usepackage{enumitem}
\usepackage[font=small,labelfont=bf]{caption}
\usepackage[colorlinks=true,linkcolor=DeepBlue,citecolor=DeepBlue,urlcolor=Cyan]{hyperref}
\usepackage{titlesec}

\definecolor{DeepBlue}{HTML}{000000}
\definecolor{Cyan}{HTML}{087F9A}
\definecolor{SoftBlue}{HTML}{EAF5F7}
\definecolor{SoftRule}{HTML}{9CC8D3}
\definecolor{Body}{HTML}{000000}

\color{Body}
\setlength{\parindent}{0pt}
\setlength{\parskip}{3.2pt}
\setlength{\abovedisplayskip}{5pt}
\setlength{\belowdisplayskip}{5pt}
\setlist{nosep,leftmargin=1.25em}
\titleformat{\section}{\large\bfseries\color{DeepBlue}}{}{0pt}{}[\vspace{-2pt}\color{Cyan}\titlerule]
\titleformat{\subsection}{\normalsize\bfseries\color{DeepBlue}}{}{0pt}{}
\titlespacing*{\section}{0pt}{7pt}{3pt}
\titlespacing*{\subsection}{0pt}{5pt}{2pt}

\newtheorem{theorem}{Theorem}
\newtheorem{corollary}[theorem]{Corollary}
\newcommand{\TV}{\operatorname{TV}}
\newcommand{\E}{\mathbb E}

\begin{document}

\begin{center}
{\LARGE\bfseries\color{DeepBlue}Finite-Horizon Indistinguishability\par
\vspace{1pt}
and the Limits of Extrapolation\par}
\vspace{4pt}
{\normalsize A two-point learning bound from hidden dynamics}\par
\vspace{6pt}
{\normalsize\textbf{Jonathan R. Landers}}\par
{\small September 2026}
\end{center}

\begin{abstract}
\noindent A learner may fit every observation in a finite experiment while remaining unable to predict a consequential latent mode outside that experiment. We formalize this as a two-point minimax bound: if two dynamical models induce training-record laws within total variation $\varepsilon$ but their future targets are separated by $\Delta$, every learner has probability at least $(1-\varepsilon)/2$ of making error at least $\Delta/2$ on one of them. A pair of weakly coupled harmonic oscillators then gives an explicit stable, bounded-energy witness. For passive noisy observation, reliable prediction of a full-amplitude consequence at time $T_*$ requires a sample count of order $(T_*/T)^4$, where $T$ is the training horizon. The statistical device is classical; the point is its interpretation as a finite-horizon limit on learning under physically incomplete model classes.
\end{abstract}

\section{Fit is not closure}

Ordinary generalization asks whether a learned rule performs on new draws governed by the observed regime. Extrapolation in a dynamical system asks a different question: could an unobserved degree of freedom be quiet during training yet control the state later? A small training error cannot answer this unless the model class already excludes such extensions.

Let an observation protocol $\mathcal O$ specify the measured variables, time window, sample budget, noise, and any allowed interventions. A physical model $M$ induces a probability law $P_M^{\mathcal O}$ for the complete record $Z$. At a later horizon $H$, let $\theta_H(M)$ be the prediction target in a metric space $(\mathcal Y,d)$. A learner is any measurable map $\widehat\theta(Z)$ from the record to $\mathcal Y$.

\begin{theorem}[Finite-horizon indistinguishability]
Suppose two models $M_0,M_1$ satisfy
\[
\TV(P_0^{\mathcal O},P_1^{\mathcal O})\leq\varepsilon,
\qquad
d(\theta_H(M_0),\theta_H(M_1))=\Delta,
\]
where $P_i^{\mathcal O}=P_{M_i}^{\mathcal O}$, $0\leq\varepsilon<1$, and $\Delta>0$. Then
\begin{align}
\inf_{\widehat\theta}\max_{i\in\{0,1\}}
P_i^{\mathcal O}\!\left\{
d(\widehat\theta,\theta_H(M_i))\geq\frac{\Delta}{2}
\right\}
&\geq\frac{1-\varepsilon}{2}, \label{eq:prob-bound}\\
\inf_{\widehat\theta}\max_{i\in\{0,1\}}
\E_i d(\widehat\theta,\theta_H(M_i))
&\geq\frac{\Delta}{4}(1-\varepsilon). \label{eq:risk-bound}
\end{align}
\end{theorem}

\begin{proof}
Given $\widehat\theta$, define a test $\phi(Z)$ by choosing the closer of the two targets, breaking ties arbitrarily. If $\phi\neq i$, the triangle inequality implies $d(\widehat\theta,\theta_H(M_i))\geq\Delta/2$. Every test obeys the two-point testing inequality
\[
P_0^{\mathcal O}(\phi=1)+P_1^{\mathcal O}(\phi=0)
\geq1-\TV(P_0^{\mathcal O},P_1^{\mathcal O}).
\]
At least one of its two error probabilities is therefore at least $(1-\varepsilon)/2$, which proves \eqref{eq:prob-bound}. Multiplying the corresponding large-error event by $\Delta/2$ proves \eqref{eq:risk-bound}.
\end{proof}

The theorem is a direct form of Le Cam's two-point method. Its learning-theoretic content is the separation of two quantities: similarity of the available training experiments and disagreement of the future targets. A learner cannot repair the first by architecture, optimization, or regularization when the second remains large.

\newpage

\section{A bounded-energy dynamical witness}

The abstract pair can be realized by stable linear mechanics. The reference model is one visible harmonic oscillator,
\[
H_0(q,p)=\tfrac12p^2+\tfrac12\omega^2q^2,
\]
with frequency $\omega>0$. The extended model adds an inaccessible coordinate $Q$ and momentum $P$:
\[
H_\delta=\tfrac12(p^2+P^2)
+\tfrac12(\omega^2+\delta^2)(q^2+Q^2)+2\omega\delta qQ,
\qquad 0<\delta\leq\omega/2.
\]
For initial state $q(0)=A>0$ and $p(0)=Q(0)=P(0)=0$, the isolated and extended visible trajectories are
\[
q_0(t)=A\cos(\omega t),
\qquad
q_\delta(t)=A\cos(\omega t)\cos(\delta t).
\]
The extension diagonalizes into normal-mode frequencies $\omega\pm\delta$. It is positive definite, conserves energy $E_\delta=\tfrac12A^2(\omega^2+\delta^2)$, and introduces a slow envelope with timescale $1/\delta$ rather than an instability.

Suppose the learner receives at most $n$ adaptively timed readings of $q$ on $[0,T]$, each corrupted by independent Gaussian noise of standard deviation $\sigma>0$. The induced record laws satisfy \cite{landers2026}
\begin{equation}
\TV(P_\delta^{\mathcal O},P_0^{\mathcal O})
\leq\min\!\left\{1,
\frac{\sqrt n\,A\delta^2T^2}{4\sigma}\right\}.
\label{eq:osc-tv}
\end{equation}
For an integer $m\geq1$, choose
\[
\delta=\frac{\omega}{4m},
\qquad
T_*=\frac{\pi}{2\delta}=\frac{2\pi m}{\omega}.
\]
Then $q_0(T_*)=A$ and $q_\delta(T_*)=0$. Shrinking $\delta$ makes the early record arbitrarily hard to distinguish while preserving a full-amplitude future separation.

\begin{corollary}[Extrapolation lower bound for the oscillator pair]
If the right-hand side of \eqref{eq:osc-tv} is at most $\varepsilon<1$, every estimator of the later visible position satisfies
\[
\max_{i\in\{0,\delta\}}
P_i^{\mathcal O}\!\left\{
|\widehat q(T_*)-q_i(T_*)|\geq A/2
\right\}
\geq\frac{1-\varepsilon}{2},
\]
and its worst-case expected absolute error is at least $A(1-\varepsilon)/4$.
\end{corollary}

This yields a concrete sample-complexity obstruction. If both models must have probability at most $\alpha<1/2$ of an error of at least $A/2$, then the theorem requires $\TV(P_\delta^{\mathcal O},P_0^{\mathcal O})\geq1-2\alpha$. Combining this necessity with \eqref{eq:osc-tv} gives
\begin{equation}
n\geq
\frac{16\sigma^2(1-2\alpha)^2}{A^2\delta^4T^4}
=\frac{256\sigma^2(1-2\alpha)^2}{\pi^4A^2}
\left(\frac{T_*}{T}\right)^4.
\label{eq:sample-complexity}
\end{equation}
Thus this family has a quartic lower bound in the ratio of prediction horizon to observation horizon. Equation \eqref{eq:sample-complexity} is necessary, not sufficient: it prices the information required to separate this particular pair and does not promise that any procedure attaining the bound exists.

Bounded active probing changes constants but not the mechanism. If a force on $q$ has total impulse at most $J$, the paper proves
\[
\TV(P_\delta^{\mathcal O},P_0^{\mathcal O})
\leq \frac{\sqrt n\,\delta^2}{2\sigma}
\left(\frac{AT^2}{2}+\frac{JT^3}{6}\right).
\]
The persistent $\delta^2$ factor means that no fixed finite probing budget gives a uniform certificate of closure over arbitrarily weak nonzero couplings.

\newpage

\section{An operational learning objective}

Absolute physical completeness is not statistically certifiable without assumptions on the allowable extensions. A useful replacement is a \emph{relative closure certificate}. Given an accepted model $M_0$, protocol $\mathcal O$, and resolution $\varepsilon$, define its observational uncertainty set
\[
\mathcal U_\varepsilon(M_0,\mathcal O)
=\{M:\TV(P_M^{\mathcal O},P_{M_0}^{\mathcal O})\leq\varepsilon\}.
\]
For future horizon $H$, define closure vulnerability
\[
V_\varepsilon(H)
=\sup_{M\in\mathcal U_\varepsilon(M_0,\mathcal O)}
d(\theta_H(M),\theta_H(M_0)).
\]
A statement of the form $V_\varepsilon(H)\leq\gamma$ says exactly what has been certified: among the specified extensions that the actual experiment resolves only to tolerance $\varepsilon$, none changes the target by more than the consequential tolerance $\gamma$ at horizon $H$. The corresponding safe horizon is
\[
H_{\mathrm{safe}}
=\sup\{H:V_\varepsilon(H)\leq\gamma\}.
\]

This suggests a practical two-gate evaluation for learned dynamical models:
\begin{enumerate}
\item \textbf{Predictive-fit gate.} Test accuracy and calibration on held-out records drawn from the measured regime.
\item \textbf{Closure-stress gate.} Search a declared extension class for latent models that remain within observational tolerance but cross the future consequence threshold.
\end{enumerate}
Failure of the second gate can be answered constructively: shorten the certified horizon, measure another variable, extend the observation window, improve resolution, or design an intervention that increases separation of the record laws.

\subsection{What the bound does and does not say}

The lower bound is uniform over a model class containing both witnesses. It does not say that a fixed nonzero coupling is forever undetectable; longer observation or richer measurement can expose it. Nor does it invalidate ordinary generalization guarantees. Those guarantees concern performance under a specified data-generating distribution, whereas the present target changes across physical models that are nearly identical on the training experiment.

Distributional diagnostics can still reduce the extension class. For example, skewness and kurtosis may reject a single-population trajectory model and force a mixture description. Such rejection establishes missing \emph{statistical structure} in the observed population. The theorem here addresses the complementary limitation: even after every available diagnostic passes, nearly identical finite-record laws may conceal dynamics that separate later. Statistical misspecification can sometimes be detected; physical closure cannot be obtained from fit alone.

The two-point inequality itself is classical \cite{lecam1986,yu1997}. The contribution of this note is therefore not a new minimax technique. It is the explicit dynamical reading of that technique, together with a stable Hamiltonian witness and a horizon-dependent sample lower bound. The formulation turns ``a good finite fit is not physical completeness'' into a testable model-risk question: how far can unresolved extensions move the target, and how much observation is required to rule them out within a declared class?

\section{Conclusion}

A finite experiment can support a genuine predictive certificate without supporting a closure claim. When two models are $\varepsilon$-close on the experiment and $\Delta$-separated at the target horizon, minimax extrapolation error remains of order $\Delta(1-\varepsilon)$. For weakly coupled oscillators, a full-amplitude future difference can coexist with arbitrarily similar early records, and the required sample count grows at least quartically with $T_*/T$. Operationally, learned dynamics should therefore be reported with both an empirical fit range and a closure-stress range relative to explicit latent extensions.

\begin{thebibliography}{9}\small
\setlength{\itemsep}{1pt}
\bibitem{landers2026}
J.~R.~Landers, ``Finite Observation and Broader Physical Processes,'' research note, 2026.

\bibitem{lecam1986}
L.~Le Cam, \emph{Asymptotic Methods in Statistical Decision Theory}. Springer, 1986.

\bibitem{yu1997}
B.~Yu, ``Assouad, Fano, and Le Cam,'' in \emph{Festschrift for Lucien Le Cam}, Springer, pp.~423--435, 1997.
\end{thebibliography}

\end{document}
